% ====================================================== decision strategy ==
\section{Decision-making strategy}

\subsection{Triage: an explainable, monotone, versioned model}

Each incident receives a priority score in $[0,112]$ from a weighted additive
model over five normalised terms plus a bounded aging term:

\[
P(i,t) \;=\; 100\sum_{k\in K} w_k\,\phi_k(i,t)
\;+\; A_{\max}\Bigl(1-e^{-\,(t-t_i)/\tau}\Bigr),
\qquad
K=\{\text{sev},\text{scale},\text{urg},\text{vuln},\text{haz}\}
\]

\noindent with $\sum_k w_k = 0.94$, $A_{\max}=12$, $\tau=900$\,s, and

\begin{center}\footnotesize
\begin{tabular}{@{}llp{72mm}@{}}
\toprule
\textbf{Term} & \textbf{Weight} & \textbf{Shape and rationale} \\
\midrule
Severity & 0.32 & $((s-1)/4)^{0.85}$, convex: catastrophic must separate sharply from serious \\
Scale & 0.18 & $\log(1+n)/\log(2001)$: 500 casualties are worse than 5 but not $100\times$ worse; saturates so mass-casualty events cannot make every other term irrelevant \\
Urgency & 0.24 & Fraction of the survivability window consumed, saturating at 1.0 past deadline so an expired window cannot monopolise the national fleet \\
Vulnerability & 0.10 & Children, elderly, disabled, isolated populations \\
Hazard & 0.10 & Environmental risk class (flood, cyclone, collapse, chemical) \\
Aging & $\le 12$ pts & Anti-starvation. \emph{Bounded}, because an unbounded aging term eventually lets a backlog of minor calls outrank a fresh catastrophe \\
\bottomrule
\end{tabular}
\end{center}

\begin{keybox}
\textbf{Why not a learned ranker at the top level?} Three reasons, all
operational rather than technical.

\begin{enumerate}
\item \emph{Explainability.} An additive model yields a per-term contribution
      breakdown for free. A gradient-boosted model yields a number and a
      SHAP plot nobody can act on at 3\,a.m.
\item \emph{Monotonicity is a safety property.} Increasing severity must never
      lower priority. Here that is true by construction and unit-tested. In a
      learned model it is testable only statistically, and a violation is a
      headline.
\item \emph{Auditability.} Every score is stamped with a policy version, so a
      post-incident inquiry replays the exact function that ran that night.
\end{enumerate}

Machine learning still belongs in this system \textemdash{} but \emph{upstream},
as feature estimators (predicted casualty count from call transcript,
predicted escalation probability, predicted service duration) feeding this
transparent aggregator. Never as the aggregator.
\end{keybox}

\subsection{The two-tier engine}

\begin{center}
\begin{tikzpicture}[node distance=5mm]
\node[svc, minimum width=30mm] (in) {Incident enriched};
\node[svc, right=8mm of in, minimum width=34mm, fill=soft!60] (t1)
  {\textbf{Tier 1 \textemdash{} Admission}\\\tiny $k$-nearest feasible, greedy\\\tiny target p99 $<$ 200\,ms};
\node[svc, right=8mm of t1, minimum width=30mm] (com) {Lease + commit\\\tiny state \texttt{PROPOSED}};
\node[svc, below=10mm of t1, minimum width=44mm, fill=soft2]
  (t2) {\textbf{Tier 2 \textemdash{} Re-optimisation} (per shard)\\\tiny exact rectangular assignment + LNS\\\tiny budgeted, anytime, debounced};
\node[ext, left=8mm of t2, minimum width=26mm] (trig)
  {Triggers:\\\tiny new incident, unit freed,\\\tiny road closed, vehicle failed,\\\tiny periodic epoch};
\draw[fl] (in) -- (t1);
\draw[fl] (t1) -- (com);
\draw[fl] (trig) -- (t2);
\draw[fl] (t2) -- ++(38mm,0) |- (com.south);
\draw[fl, dashed] (com.south) -- ++(0,-5mm) -| node[lbl, pos=0.25, above] {upgrade if gain $>$ hysteresis} (t2.east);
\end{tikzpicture}
\end{center}

Both tiers evaluate the \emph{same} cost function. This matters more than it
looks: if the fast path and the optimal path disagreed about units, every
re-optimisation epoch would churn simply because the two tiers were
speaking different languages.

% ======================================================== optimisation =====
\section{Optimization approach}

\subsection{The model}

Let $R$ be the set of candidate resources and $I$ the set of open incidents.
Incident $i$ requires $q_{i\rho}$ units of resource type $\rho$. The natural
formulation is

\begin{align*}
\min_{x} \quad & \sum_{r\in R}\sum_{i\in I} c_{ri}\,x_{ri} \;+\; \sum_{i} \Pi_i u_i \\
\text{s.t.}\quad
& \textstyle\sum_{r\in R_\rho} x_{ri} + u_{i\rho} = q_{i\rho} && \forall i,\rho \quad\text{(demand)}\\
& \textstyle\sum_{i\in I} x_{ri} \le 1 && \forall r \quad\text{(a unit serves one incident)}\\
& \textstyle\sum_{i} x_{ri}\,\beta_i \le \text{cap}_h && \forall h \quad\text{(hospital beds)}\\
& x_{ri}\in\{0,1\}
\end{align*}

\noindent which is a \emph{generalized assignment problem with side
constraints} \textemdash{} NP-hard. The engineering question is not "which
solver", it is "which reformulation makes the dominant sub-problem exactly
solvable inside the time budget".

\begin{keybox}
\textbf{Requirement expansion.} An incident needing 2 ambulances and 1 rescue
team becomes three independent unit-\emph{slots}. After expansion the demand
constraint becomes "each slot gets at most one unit", and the core problem is
a \textbf{rectangular linear assignment problem}: solvable \emph{exactly} in
$O(n^3)$ by the Jonker--Volgenant shortest-augmenting-path algorithm.

The genuinely combinatorial residue \textemdash{} bed capacity, coverage
reserve, multi-stop sequencing \textemdash{} is then handled by penalty terms
in $c_{ri}$ and by a Large Neighbourhood Search pass over the exact solution,
rather than by abandoning optimality altogether.
\end{keybox}

\subsection{The cost function}

For resource $r$ and slot $s$ belonging to incident $i$, in \emph{urgency-minutes}:

\[
c_{rs} =
\underbrace{\mu_i\,(\eta_{rs} - H)}_{\text{travel}}
+ \underbrace{\lambda\,\mu_i\,[\eta_{rs}-\sigma_i]^+}_{\text{deadline}}
+ \underbrace{\omega\,[\kappa_r-\kappa_s]^+}_{\text{over-qualification}}
+ \underbrace{\gamma\,S_\rho}_{\text{scarcity}}
+ \underbrace{\theta\,C_r}_{\text{coverage}}
+ \underbrace{\mu_i\,\Sigma_r}_{\text{stability}}
\]

\noindent where $\mu_i = \max(0.15, P_i/100)$ is normalised urgency,
$\eta_{rs}$ is estimated door-to-scene minutes, $H$ the planning horizon,
$\sigma_i$ the remaining deadline slack, $\kappa$ capability levels, $S_\rho$
the scarcity of the asset class, $C_r$ the coverage penalty, and $\Sigma_r$
the stability term.

\begin{center}\footnotesize
\begin{tabular}{@{}lp{104mm}@{}}
\toprule
\textbf{Term} & \textbf{What it encodes and why it is needed} \\
\midrule
Travel & Weighted response time. Negative for useful assignments, so the solver is drawn toward filling urgent slots with near units. \\
Deadline & Convex penalty for arriving after the survivability window; zero if the unit makes it in time. \\
Over-qualification & Discourages spending a level-3 advanced unit on a slot that a level-1 unit satisfies, preserving capability for cases that need it. \\
Scarcity & A helicopter spent on a routine transfer is a helicopter unavailable for the case only it can reach. Deliberately \emph{not} scaled by urgency, so the disincentive is strongest exactly where the asset is least warranted. \\
Coverage & Marginal cost of stripping a district of its last reserve, weighted by that district's forecast arrival rate. Charged only on draws from the idle pool. \\
Stability & Continuity bonus for staying on the current task; switching penalty scaling with assignment maturity. \\
\bottomrule
\end{tabular}
\end{center}

\begin{decbox}
\dec{Hard reach radius.}{A unit that cannot reach the scene within
$\eta_{\max}=100$ minutes is infeasible at any price, not merely expensive.
This is not a tuning constant, it is a correctness guard: a minimum-cost
assignment algorithm assigns \emph{every row it is given}, so without an
explicit feasibility cutoff an exact solver will cheerfully dispatch an
ambulance across the country, remove it from the pool for six hours, and
still arrive long after the clinical window has closed. This was observed and
fixed during development of the reference implementation.}
\end{decbox}

\subsection{Churn control: stability as part of the objective}

Re-optimisation is only valuable if it does not thrash crews. Three mechanisms
work together:

\begin{enumerate}
\item \textbf{Assignment maturity.} \texttt{PROPOSED} units are free to move;
      \texttt{COMMITTED} units cost 40 urgency-minutes to move;
      \texttt{EN\_ROUTE} units cost 160; \texttt{ARRIVED} units are
      irrevocable. Optimality is thereby \emph{defined} to include the cost of
      changing one's mind.
\item \textbf{Progress-aware estimates.} A unit already travelling is costed on
      its \emph{remaining} time from its \emph{projected current position}, not
      re-quoted from the depot it left twenty minutes ago. Without this the
      solver perpetually believes a moving vehicle is still at its origin and
      swaps it out endlessly.
\item \textbf{Demand-level supersession.} A unit is stood down only when its
      \emph{(incident, resource type)} demand is fully staffed by others.
      Comparing individual slot indices instead reports a unit as displaced
      whenever the solver merely permutes interchangeable columns.
\end{enumerate}

\begin{keybox}
These three are not polish. Before they were implemented the reference
engine issued 2{,}233 assignments for 148 incidents \textemdash{} a
15$\times$ churn rate that made outcomes \emph{worse} than the naive
baseline despite solving the assignment problem optimally. Afterwards churn
falls below 2\% of commits and the optimiser's advantage appears. The lesson
generalises: in a continuously re-planned system, plan stability is not a
usability concern, it is a prerequisite for the optimisation to deliver any
value at all.
\end{keybox}

\subsection{Solver ladder}

\begin{center}\footnotesize
\begin{tabular}{@{}llp{62mm}@{}}
\toprule
\textbf{Stage} & \textbf{Complexity} & \textbf{Role} \\
\midrule
Greedy over sorted cost & $O(nm\log nm)$ & Warm start; also the Tier-1 admission decision \\
Jonker--Volgenant & $O(n^3)$ & Exact optimum of the expanded linear problem \\
LNS (ruin \& recreate) & anytime & Improves under coupling constraints the linear model cannot express; never returns worse than its input \\
\bottomrule
\end{tabular}
\end{center}

\noindent The exact solver is verified against exhaustive enumeration on small
instances and asserted never to be worse than greedy on random instances. LNS
is asserted to be monotonically non-worsening. All three properties are
unit tests, not claims.

\begin{decbox}
\dec{Why not a MILP solver over the full model?}{A national fleet gives
per-shard matrices of $10^2$--$10^3$, where Jonker--Volgenant runs in
single-digit milliseconds. A MILP over the unexpanded model with capacity
constraints does not reliably fit a sub-second budget and, critically, has no
useful \emph{anytime} behaviour: a truncated branch-and-bound may hold no
feasible incumbent at all. AEGIS always has a shippable plan.}
\end{decbox}

% ================================================== continuous re-opt ======
\section{Continuous re-optimization}

The environment changes in five ways, each with a defined response:

\begin{center}\footnotesize
\begin{tabular}{@{}lp{46mm}p{52mm}@{}}
\toprule
\textbf{Change} & \textbf{Detection} & \textbf{Response} \\
\midrule
New incident & \texttt{incidents.enriched} & Tier-1 admission immediately; shard queued for re-optimisation \\
Unit freed & Telemetry state transition & Debounced re-optimisation trigger; unit re-enters the pool \\
Road closed / flooded & \texttt{environment.events} & Surgical cache invalidation for affected cells only; shard re-optimised \\
Vehicle failure & Telemetry heartbeat loss or explicit report & Live assignment revoked with reason code, incident returns to the unmet pool, re-optimisation triggered \\
Hospital saturation & Bed reservation CAS failure & Handover diverted to next-nearest facility with capacity; if none, the unit holds on scene and the shortage propagates visibly into fleet availability \\
\bottomrule
\end{tabular}
\end{center}

\begin{keybox}
\textbf{Debouncing.} Running the solver on every environment change would be
both wasteful and destabilising \textemdash{} it is the standard control-loop
problem of a noisy input signal. Triggers are coalesced into at most one epoch
per debounce window (15\,s in the reference configuration), with a periodic
floor epoch (30\,s) so that a quiet system still refreshes its plan.
\end{keybox}

% =================================================== conflict prevention ===
\section{Preventing resource conflicts}
\label{sec:conflict}

Double-booking an ambulance is the failure mode that ends the programme. It is
guarded at three independent layers, so that no single mechanism is
load-bearing alone.

\begin{center}\footnotesize
\begin{tabular}{@{}lp{44mm}p{50mm}@{}}
\toprule
\textbf{Layer} & \textbf{Mechanism} & \textbf{What it protects against} \\
\midrule
Partitioning & Single writer per region shard & The common case: two workers planning the same fleet \\
Reservation & TTL lease in Redis, acquired before commit, released with a fencing token & Cross-region borrowing; a stale holder freeing somebody else's lease \\
Persistence & Optimistic concurrency \\ \texttt{UPDATE \dots WHERE version = ?} & Lost updates when two commits interleave at the database \\
\bottomrule
\end{tabular}
\end{center}

\subsection{Why leases have a TTL}

If the optimiser crashes between reserving a unit and committing the
assignment, a lock without expiry strands that ambulance \emph{forever}. It is
a correctness bug that is invisible in testing and catastrophic in production,
because it manifests as a slow, silent erosion of fleet capacity. A TTL lease
plus a background reaper makes crash recovery automatic: the resource returns
to the pool without operator intervention. The reference implementation counts
reclaimed leases as a first-class metric.

\subsection{Group reservation is all-or-nothing}

A collapsed building requiring 2 ambulances and 1 rescue team must not end up
holding a rescue team hostage while the ambulances go elsewhere. Partial
allocation both wastes the committed unit and \emph{hides the shortfall} from
the operator, who sees an incident marked "assigned". Group reservation is
therefore a small saga: acquire all leases, and on any failure compensate by
releasing everything acquired so far.

\subsection{Optimistic rather than pessimistic concurrency}

Dispatch conflicts are rare \textemdash{} two shards genuinely contending for
one unit is an edge case \textemdash{} but lock convoys under surge are fatal.
Optimistic concurrency pays nothing in the common case and forces a re-read
and retry in the rare one. Pessimistic row locks invert that trade exactly
when the system is busiest.

\begin{keybox}
\textbf{Verification.} The invariant is asserted, not assumed. A
multi-threaded contention test races twelve workers over twenty-five units for
2{,}400 attempts and asserts that every unit has exactly one holder and that
contention actually occurred. The full simulation additionally runs an
integrity audit over the final state: zero conflicts in every configuration
and every seed.
\end{keybox}
